% GENERATED FILE. Edit canonical lessons, then run npm run generate:books.
% canonical-source-hash: fnv1a64:56626790d9817e35
% canonical-lessons: GU-C276-khaine, GU-C276-chevte, GU-C276-pahela-pachi, GU-C276-ek-ravivar, GU-R276-a-story-in-order

\chapter{A story in order}
\label{ch:gu-a-story-in-order}
\hlchaptermodality{276}

\begin{chapteropening}
\textbf{By the end of this chapter:} \emph{I can join two actions with the -īne form (jaīne, khāīne) and -yā pachī, keep a story in order with pahelā̃, pachī, tyār pachī and chevṭe, and read a short story told in the past.}

Complete GU-R276-a-story-in-order: tell a short story in order, and read one.
\end{chapteropening}

\section[\texorpdfstring{\gu{ખાઈને} · \gu{જમ્યા} \gu{પછી} (khāīne · jamyā pachī)}{ખાઈને · જમ્યા પછી (khāīne · jamyā pachī)}]{\gu{ખાઈને}, \gu{જમ્યા} \gu{પછી} (khāīne, jamyā pachī) — having done, after doing}

\label{lesson:GU-C276-khaine}



\begin{quote}
Say \textquotedblleft{}please\textquotedblright{}, the long way. (\textbf{\gu{મહેરબાની} \gu{કરીને}}.) Its \textbf{\gu{કરીને}} is the form this lesson teaches. Then say \textquotedblleft{}I went home\textquotedblright{}. (\emph{hũ ghare gayo}.)
\end{quote}

\begin{grammarlens}[title={Grammar Lens: the -\gu{ઈને} form}]
\begin{quote}
\textbf{\gu{હું} \gu{ઘરે} \gu{જઈને} \gu{જમ્યો}.}
\end{quote}

\begin{quote}
\emph{hũ ghare jaīne jamyo.} — \textquotedblleft{}I went home and ate.\textquotedblright{}
\end{quote}

One doer, two actions in a row: the first takes \textbf{-\gu{ઈને}}, \textquotedblleft{}having done it\textquotedblright{}, and the second ends the sentence. After a vowel the stem takes \textbf{-\gu{ઈને}}: \textbf{\gu{જઈને}}, \textbf{\gu{ખાઈને}}, \textbf{\gu{લઈને}}, \textbf{\gu{જોઈને}}. After a consonant it takes \textbf{-\gu{ીને}}: \textbf{\gu{કરીને}}, \textbf{\gu{લખીને}}, \textbf{\gu{વાંચીને}}.

The last verb decides the doer's form: \textbf{\gu{હું} \gu{ખાઈને} \gu{આવ્યો}} (\emph{hũ khāīne āvyo}), \textquotedblleft{}I ate and came\textquotedblright{}, but \textbf{\gu{મેં} \gu{ખાઈને} \gu{ચા} \gu{પીધી}} (\emph{mẽ khāīne cā pīdhī}), \textquotedblleft{}I ate and then drank tea\textquotedblright{}.
\end{grammarlens}

\subsection*{You'll want to know: -\gu{્યા} \gu{પછી}}
\begin{quote}
\textbf{\gu{જમ્યા} \gu{પછી} \gu{અમે} \gu{બજાર} \gu{ગયા}.}
\end{quote}

\begin{quote}
\emph{jamyā pachī ame bajār gayā.} — \textquotedblleft{}After eating, we went to the market.\textquotedblright{}
\end{quote}

\textbf{\gu{પછી}} after the past stem in \textbf{-\gu{્યા}} says \textquotedblleft{}after doing it\textquotedblright{}: \textbf{\gu{જમ્યા} \gu{પછી}}, \textbf{\gu{કામ} \gu{કર્યા} \gu{પછી}} (\emph{kām karyā pachī}), \textbf{\gu{ચા} \gu{પીધા} \gu{પછી}} (\emph{cā pīdhā pachī}). That stem takes no \textbf{-\gu{એ}} doer of its own: \textbf{\gu{કામ} \gu{કર્યા} \gu{પછી} \gu{હું} \gu{ઘરે} \gu{ગયો}}.

\subsection*{Your turn}
Say these aloud:

\begin{itemize}
\raggedright
  \item \textquotedblleft{}Mira wrote the letter and sent it\textquotedblright{} — \emph{Mīrāe patra lakhīne mokalyo}
  \item \textquotedblleft{}After the work, I drank tea\textquotedblright{} — \emph{kām karyā pachī mẽ cā pīdhī}
\end{itemize}

\begin{tcolorbox}[breakable,colback=teal!4,colframe=teal!35!black,title={Before you move on}]
\textquotedblleft{}Having gone\textquotedblright{}? (\textbf{\gu{જઈને}}.) \textquotedblleft{}Having done\textquotedblright{}? (\textbf{\gu{કરીને}}.) \textquotedblleft{}After eating\textquotedblright{}? (\textbf{\gu{જમ્યા} \gu{પછી}}.) In \textbf{\gu{મેં} \gu{ખાઈને} \gu{ચા} \gu{પીધી}}, which verb chose \textbf{\gu{મેં}}? (\textbf{\gu{પીધી}}.)
\end{tcolorbox}
\section[\texorpdfstring{\gu{છેવટે} (chevṭe)}{છેવટે (chevṭe)}]{\gu{છેવટે} (chevṭe) — finally; in the end}

\label{lesson:GU-C276-chevte}



\begin{quote}
Before the new one: say \textquotedblleft{}I went home and ate\textquotedblright{}. (\emph{hũ ghare jaīne jamyo}.) Then the word for \textquotedblleft{}last\textquotedblright{}. (\textbf{\gu{છેલ્લું}}, \emph{chellũ}.)
\end{quote}

\subsection*{You'll want to know: \gu{છેવટે}}
\textbf{\gu{છેવટે}} — \emph{chevṭe} — \textquotedblleft{}finally; in the end\textquotedblright{}.

\begin{quote}
\textbf{\gu{છેવટે} \gu{ટ્રેન} \gu{આવી}.}
\end{quote}

\begin{quote}
\emph{chevṭe ṭren āvī.} — \textquotedblleft{}In the end the train came.\textquotedblright{}
\end{quote}

\textbf{\gu{છેવટે}} comes after waiting, or after a long run of things: it closes them off.

\begin{quote}
\textbf{\gu{અમે} \gu{આખો} \gu{દિવસ} \gu{ચાલ્યા}, \gu{અને} \gu{છેવટે} \gu{ગામ} \gu{પહોંચ્યા}.}
\end{quote}

\begin{quote}
\emph{ame ākho divas cālyā, ane chevṭe gām pahõcyā.} — \textquotedblleft{}We walked all day, and finally reached the village.\textquotedblright{}
\end{quote}

\begin{cousinweb}
\textbf{\gu{છેવટે}} (\emph{chevṭe}) is built from \textbf{\gu{છેવટ}}, \textquotedblleft{}the end\textquotedblright{}, and the \textbf{-\gu{ે}} of \textbf{\gu{સવારે}} and \textbf{\gu{ઘરે}}: \textquotedblleft{}at the end\textquotedblright{}. CDIAL lists Gujarati \emph{chevaṭ}, \textquotedblleft{}end\textquotedblright{}, under Sanskrit \emph{cheda}, \textquotedblleft{}a cut, a break\textquotedblright{} (CDIAL 5064), the entry that also holds the \emph{che-} of \textbf{\gu{છેલ્લું}}, \textquotedblleft{}last\textquotedblright{}.
\end{cousinweb}

\subsection*{Your turn}
Say these aloud:

\begin{itemize}
\raggedright
  \item \emph{chevṭe}
  \item \textquotedblleft{}In the end the train came\textquotedblright{} — \emph{chevṭe ṭren āvī}
  \item \textquotedblleft{}Finally we got home\textquotedblright{} — \emph{chevṭe ame ghare pahõcyā}
\end{itemize}

\begin{tcolorbox}[breakable,colback=teal!4,colframe=teal!35!black,title={Before you move on}]
What does \textbf{\gu{છેવટે}} mean? (\textquotedblleft{}Finally; in the end\textquotedblright{}.) Which word that you know shares its start? (\textbf{\gu{છેલ્લું}}, \textquotedblleft{}last\textquotedblright{}.)
\end{tcolorbox}
\section[\texorpdfstring{\gu{પહેલાં} · \gu{પછી} … (pahelā̃ · pachī …)}{પહેલાં · પછી … (pahelā̃ · pachī …)}]{\gu{પહેલાં}, \gu{પછી}, \gu{ત્યાર} \gu{પછી} (pahelā̃, pachī, tyār pachī) — first, then, after that}

\label{lesson:GU-C276-pahela-pachi}



\begin{quote}
Say \textquotedblleft{}before\textquotedblright{} and \textquotedblleft{}later\textquotedblright{}. (\textbf{\gu{પહેલાં}}, \textbf{\gu{પછી}}.) Then \textquotedblleft{}in the end the train came\textquotedblright{}. (\emph{chevṭe ṭren āvī}.)
\end{quote}

\begin{grammarlens}[title={Grammar Lens: first, then}]
\begin{quote}
\textbf{\gu{પહેલાં} \gu{અમે} \gu{ચા} \gu{પીધી}, \gu{પછી} \gu{અમે} \gu{બજાર} \gu{ગયા}.}
\end{quote}

\begin{quote}
\emph{pahelā̃ ame cā pīdhī, pachī ame bajār gayā.} — \textquotedblleft{}First we drank tea, then we went to the market.\textquotedblright{}
\end{quote}

At the head of a sentence, \textbf{\gu{પહેલાં}} is \textquotedblleft{}first\textquotedblright{} and \textbf{\gu{પછી}} is \textquotedblleft{}then\textquotedblright{}. Each sentence keeps its own past: \textbf{\gu{અમે} \gu{ચા} \gu{પીધી}} follows the tea, \textbf{\gu{અમે} \gu{બજાર} \gu{ગયા}} follows the people who went.
\end{grammarlens}

\subsection*{You'll want to know: \gu{ત્યાર} \gu{પછી} and \gu{છેવટે}}
\begin{quote}
\textbf{\gu{ત્યાર} \gu{પછી} \gu{અમે} \gu{ઘરે} \gu{આવ્યા}.}
\end{quote}

\begin{quote}
\emph{tyār pachī ame ghare āvyā.} — \textquotedblleft{}After that we came home.\textquotedblright{}
\end{quote}

\textbf{\gu{ત્યાર} \gu{પછી}} (\emph{tyār pachī}), \textquotedblleft{}after that\textquotedblright{}, is a fuller \textbf{\gu{પછી}}: it opens a new step. A story runs \textbf{\gu{પહેલાં}} … \textbf{\gu{પછી}} … \textbf{\gu{ત્યાર} \gu{પછી}} … and ends on \textbf{\gu{છેવટે}}:

\noindent
\begin{tabularx}{\linewidth}{@{}>{\raggedright\arraybackslash}X>{\raggedright\arraybackslash}X@{}}
\toprule
\textbf{step} & \textbf{word} \\
\midrule
the first thing & \textbf{\gu{પહેલાં}} \\
the next & \textbf{\gu{પછી}} \\
a new step & \textbf{\gu{ત્યાર} \gu{પછી}} \\
the last & \textbf{\gu{છેવટે}} \\
\bottomrule
\end{tabularx}

\subsection*{Your turn}
Say these aloud:

\begin{itemize}
\raggedright
  \item \textquotedblleft{}First I wrote the letter, then I drank tea\textquotedblright{} — \emph{pahelā̃ mẽ patra lakhyo, pachī mẽ cā pīdhī}
  \item \textquotedblleft{}In the end we came home in the evening\textquotedblright{} — \emph{chevṭe ame sānje ghare āvyā}
\end{itemize}

\begin{tcolorbox}[breakable,colback=teal!4,colframe=teal!35!black,title={Before you move on}]
\textquotedblleft{}First\textquotedblright{}? (\textbf{\gu{પહેલાં}}.) \textquotedblleft{}After that\textquotedblright{}? (\textbf{\gu{ત્યાર} \gu{પછી}}.) Which word ends the story? (\textbf{\gu{છેવટે}}.)
\end{tcolorbox}
\section[(reading)]{(reading) — a Sunday at the village}

\label{lesson:GU-C276-ek-ravivar}



\begin{quote}
Mira tells of a Sunday at her grandfather's village. Every word is one you have met. Listen for the order words: \textbf{\gu{પહેલાં}}, \textbf{\gu{ત્યાર} \gu{પછી}}, \textbf{\gu{છેવટે}}.
\end{quote}

\begin{tcolorbox}[breakable,skin=enhanced,colback=orange!6,colframe=orange!45!black,boxrule=0.5pt,arc=1mm,left=6pt,right=6pt,top=4pt,bottom=4pt,fonttitle=\bfseries,title={Read this}]
\begin{quote}
\textbf{\gu{ગઈકાલે} \gu{રવિવાર} \gu{હતો}. \gu{સવારે} \gu{અમે} \gu{બસમાં} \gu{ગામ} \gu{ગયા}. \gu{બસમાં} \gu{ઘણી} \gu{ભીડ} \gu{હતી}.} \textbf{\gu{પહેલાં} \gu{અમે} \gu{દાદાના} \gu{ઘરે} \gu{ગયા}. \gu{દાદા} \gu{ઓસરીમાં} \gu{બેસીને} \gu{ચા} \gu{પીતા} \gu{હતા}. \gu{દાદાએ} \gu{અમને} \gu{પણ} \gu{ચા} \gu{આપી}.} \textbf{\gu{ત્યાર} \gu{પછી} \gu{અમે} \gu{ખેતરમાં} \gu{ગયા}. \gu{ત્યાં} \gu{બળદ} \gu{ચરતા} \gu{હતા}. \gu{અમે} \gu{ખેતરમાં} \gu{ચાલતા} \gu{હતા} \gu{ત્યારે} \gu{અચાનક} \gu{વરસાદ} \gu{પડ્યો}.} \textbf{\gu{અમે} \gu{દોડીને} \gu{ઘરે} \gu{આવ્યા}. \gu{શું} \gu{બન્યું}? \gu{અમારાં} \gu{કપડાં} \gu{ભીનાં} \gu{થઈ} \gu{ગયાં} \gu{હતાં}!} \textbf{\gu{જમ્યા} \gu{પછી} \gu{દાદાએ} \gu{અમને} \gu{એક} \gu{જૂની} \gu{વાર્તા} \gu{કહી}.} \textbf{\gu{સાંજે} \gu{અમે} \gu{બસમાં} \gu{શહેર} \gu{પાછા} \gu{જવાના} \gu{હતા}, \gu{પણ} \gu{બસ} \gu{મોડી} \gu{હતી}. \gu{છેવટે} \gu{અમે} \gu{રાતે} \gu{ઘરે} \gu{પહોંચ્યા}.} \textbf{\gu{દિવસ} \gu{લાંબો} \gu{હતો}, \gu{પણ} \gu{મને} \gu{મજા} \gu{આવી}. \gu{આવતા} \gu{રવિવારે} \gu{હું} \gu{ગામ} \gu{નહીં} \gu{જાઉં}. \gu{દાદા} \gu{શહેર} \gu{આવશે}.}
\end{quote}

\emph{Read it:} the story once, straight through, without translating

\emph{Read it:} it again, and sort the verbs: scene, event or plan
\end{tcolorbox}

\begin{culture}
New pieces: \textbf{\gu{દાદાના} \gu{ઘરે}}, \textquotedblleft{}at Grandfather's house\textquotedblright{}; \textbf{\gu{બેસીને}} (\emph{besīne}), \textquotedblleft{}sitting\textquotedblright{}; \textbf{\gu{ભીનાં} \gu{થઈ} \gu{ગયાં} \gu{હતાં}}, \textquotedblleft{}had got wet\textquotedblright{}, neuter plural for \textbf{\gu{કપડાં}}.

The scenes take \textbf{-\gu{તા} \gu{હતા}}; the events take the plain past. \textbf{\gu{દાદાએ}} has \textbf{-\gu{એ}} because he gave and told something: \textbf{\gu{ચા} \gu{આપી}}, \textbf{\gu{વાર્તા} \gu{કહી}}. The bus was a plan that changed: \textbf{\gu{જવાના} \gu{હતા}, \gu{પણ} …}. \textbf{\gu{નહીં} \gu{જાઉં}} is the short form of \textbf{\gu{જવું}} after \textbf{\gu{નહીં}}.
\end{culture}

\subsection*{Your turn}
Say these aloud:

\begin{itemize}
\raggedright
  \item what Grandfather was doing on the veranda
  \item why they got home late
\end{itemize}

\begin{tcolorbox}[breakable,colback=teal!4,colframe=teal!35!black,title={Before you move on}]
Which word turns the story? (\textbf{\gu{અચાનક}}.) Why \textbf{\gu{ભીનાં} \gu{થઈ} \gu{ગયાં} \gu{હતાં}}, not \textbf{\gu{ભીનું} \gu{થઈ} \gu{ગયું}}? (\textbf{\gu{કપડાં}} is neuter plural.) Who is coming next Sunday? (Grandfather.)
\end{tcolorbox}
\section[(review)]{Review — a story in order}

\label{lesson:GU-R276-a-story-in-order}



\begin{quote}
Say \textquotedblleft{}I was reading\textquotedblright{}. (\emph{hũ vā̃cto hato}.) Then \textquotedblleft{}what happened?\textquotedblright{}, in a story. (\emph{śũ banyũ?})
\end{quote}

\subsection*{Your turn}
Say these aloud:

\begin{itemize}
\raggedright
  \item I went home and ate.
  \item After eating, we went to the market.
  \item First we drank tea, then we went to the market.
  \item After that we came home, and in the end we slept.
  \item In Mira's Sunday, what happened in the field?
\end{itemize}

\begin{tcolorbox}[breakable,colback=teal!4,colframe=teal!35!black,title={Before you move on}]
\textquotedblleft{}Having eaten\textquotedblright{}? (\textbf{\gu{ખાઈને}}.) \textquotedblleft{}After the work\textquotedblright{}? (\textbf{\gu{કામ} \gu{કર્યા} \gu{પછી}}.) Put these in order: \textbf{\gu{છેવટે}}, \textbf{\gu{પહેલાં}}, \textbf{\gu{ત્યાર} \gu{પછી}}, \textbf{\gu{પછી}}. (\textbf{\gu{પહેલાં}}, \textbf{\gu{પછી}}, \textbf{\gu{ત્યાર} \gu{પછી}}, \textbf{\gu{છેવટે}}.)
\end{tcolorbox}
